\documentclass[11pt,letterpaper,oneside]{report}
\usepackage[T1]{fontenc}
\usepackage[utf8]{inputenc}
\usepackage[indonesian]{babel}
\usepackage[margin=1in]{geometry}
\usepackage{amsmath}
\usepackage{newtxtext,newtxmath}
\usepackage{microtype}
\usepackage{graphicx,booktabs,tabularx,array}
\usepackage[font=small,labelfont=bf]{caption}
\usepackage{titlesec,xurl}
\usepackage[hidelinks,unicode]{hyperref}

\graphicspath{{figures/}}
\titleformat{\chapter}[display]
  {\normalfont\bfseries\Large}{\MakeUppercase{\chaptername} \thechapter}{5pt}{\huge}
\titlespacing*{\chapter}{0pt}{0pt}{20pt}
\newcolumntype{Y}{>{\raggedright\arraybackslash}X}
\newcommand{\pct}[1]{#1\%}
\newcommand{\kode}[1]{\nolinkurl{#1}}
\hypersetup{
  pdftitle={Pengenalan Ucapan Bahasa Sunda dengan wav2vec 2.0},
  pdfauthor={Nathaniel Jonathan Rusli, Maheswara Bayu Kaindra, Peter Wongsoredjo, Jessica Allen, Haegen Quinston},
  pdfsubject={Laporan IF4071 Pemrosesan Ucapan, Tugas Besar Fase 2}
}

\begin{document}
\hypersetup{pageanchor=false}
\begin{titlepage}
  \centering
  \vspace*{0.2in}
  {\large\bfseries IF4071 Pemrosesan Ucapan\par}
  \vspace{0.25in}
  {\LARGE\bfseries Pengenalan Ucapan Bahasa Sunda\\[0.15em] dengan wav2vec 2.0\par}
  \vspace{0.19in}
  {\large Laporan Tugas Besar Fase 2\par}
  \vfill
  \includegraphics[width=2.4in]{logo-itb.png}\par
  \vfill
  {\large oleh\par}
  \vspace{0.12in}
  {\large
  \begin{tabular}{@{}l@{\hspace{2.5em}}r@{}}
    Nathaniel Jonathan Rusli & 13523013 \\
    Maheswara Bayu Kaindra   & 13523015 \\
    Peter Wongsoredjo        & 13523039 \\
    Jessica Allen            & 13523059 \\
    Haegen Quinston          & 13523109
  \end{tabular}\par}
  \vfill
  {\large\bfseries
    SEKOLAH TEKNIK ELEKTRO DAN INFORMATIKA\par
    \vspace{0.1in}
    INSTITUT TEKNOLOGI BANDUNG\par
  }
  \vspace{0.16in}
  {\large Semester 1 Tahun Akademik 2026/2027\par}
\end{titlepage}

\pagenumbering{roman}
\hypersetup{pageanchor=true}
\chapter*{Abstrak}
\addcontentsline{toc}{chapter}{Abstrak}

\noindent\textit{Draf kolaborasi. Abstrak akan ditinjau setelah Bab 5--8 selesai.}\par\medskip

Tugas besar ini mengembangkan sistem pengenalan ucapan otomatis bahasa Sunda dengan menyetel ulang dua model wav2vec 2.0 XLSR-53 berbasis \textit{Connectionist Temporal Classification} (CTC). Korpus terdiri atas 2.185 ujaran dengan durasi sekitar 3 jam 32 menit dari 21 penutur. Himpunan latih, validasi, dan uji menggunakan penutur yang saling terpisah. Empat laju pembelajaran diuji, lalu bobot model dipilih berdasarkan \textit{Word Error Rate} (WER) pada data validasi. Setelah penyetelan, model Cahya memperoleh WER \pct{44,44} dan \textit{Character Error Rate} (CER) \pct{20,94} pada validasi, sedangkan model Indonesian NLP memperoleh WER \pct{41,59} dan CER \pct{20,41}. Pada 327 ujaran uji, penyetelan model Indonesian NLP menurunkan WER dari \pct{66,57} menjadi \pct{41,06} dan CER dari \pct{22,96} menjadi \pct{14,73}. Substitusi mendominasi kesalahan kata. Proporsi kata referensi yang salah lebih tinggi untuk kata yang tidak muncul dalam data latih, dan tingkat kesalahan berbeda antarpenutur. Evaluasi uji mencakup empat penutur dalam satu pembagian data.

\medskip
\noindent\textbf{Kata kunci} \enspace pengenalan ucapan otomatis, bahasa Sunda, wav2vec 2.0, CTC, WER, CER.

\clearpage
\tableofcontents
\clearpage
\pagenumbering{arabic}

\chapter{Pendahuluan}
\label{bab:pendahuluan}

\section{Latar belakang}

Pengenalan ucapan otomatis (\textit{Automatic Speech Recognition} atau ASR) mengubah rekaman ucapan menjadi teks. Jumlah rekaman bahasa Sunda bertranskripsi dalam tugas ini terbatas. Ujaran spontan juga memuat kata pengisi, pengulangan, alih kode, dan variasi pengucapan. Perbedaan antara \textit{e} dan \textit{é} perlu dipertahankan agar penilaian mencerminkan ejaan transkripsi rujukan.

Pada Fase 2, korpus kelompok dari Fase 1 digabungkan dengan korpus G04 untuk menyetel ulang model pralatih wav2vec 2.0 sesuai spesifikasi tugas IF4071 \cite{spesifikasi,fase1,korpusg04}. Representasi akustik dari model pralatih memberi titik awal ketika data lokal bertranskripsi terbatas \cite{baevski2020,conneau2021}.

\section{Tujuan}

\begin{enumerate}
  \item Membangun sistem ASR bahasa Sunda dengan menyetel ulang model wav2vec 2.0 berbasis CTC menggunakan pasangan audio dan transkripsi Fase 1.
  \item Mengukur perubahan WER dan CER terhadap kondisi \textit{zero-shot} pada penutur yang terpisah dari data latih. Evaluasi ini juga membandingkan dua titik awal model pada data validasi.
  \item Menganalisis pola kesalahan pada data uji serta merumuskan perbaikan data, evaluasi, dan dekode.
\end{enumerate}

\section{Repositori tugas besar}

Kode prapemrosesan, \textit{notebook} pelatihan dan inferensi, manifes, pembagian data, serta hasil evaluasi tersedia di folder \kode{ASR} pada cabang \kode{asr-dev} repositori GitHub publik. Sumber artefak untuk laporan ini adalah \textit{commit} \kode{fa3ebaa} \cite{repoasr}.

\begin{quote}
\url{https://github.com/0xNathaniel/sundanese-asr/tree/asr-dev/ASR}
\end{quote}

\chapter{Korpus Bahasa Sunda}
\label{bab:korpus}

\section{Sumber rekaman dan anotasi}

Korpus kelompok memuat 750 ujaran dari 10 penutur berusia 15 sampai 51 tahun. Delapan penutur jati berbicara spontan dan dua penutur lainnya membacakan naskah. Topiknya mencakup kegiatan sehari-hari, keluarga, makanan, pekerjaan atau pendidikan, serta hobi. Rekaman ponsel umumnya dibuat di dalam ruangan. Anotasi mempertahankan kata serapan dan pengulangan kata. Anotasi juga menandai kata pengisi dan kejadian khusus \cite{fase1,korpusutama}.

Korpus G04 menambah 1.435 ujaran dari 11 penutur dengan status penutur jati \textit{Unknown}. Topiknya meliputi perkenalan, kegiatan sehari-hari, keluarga, makanan, pendidikan, pekerjaan, dan perjalanan. Setiap korpus memiliki skema kode topiknya sendiri \cite{korpusg04,repoasr}.

\begin{table}[!htbp]
  \centering\small
  \caption{Komposisi korpus setelah seleksi. Durasi dihitung dari berkas audio pada manifes.}
  \label{tab:korpus}
  \begin{tabular}{@{}lrrrrr@{}}
    \toprule
    Sumber & Penutur & Ujaran & Durasi (menit) & \textit{Good} & \textit{Acceptable} \\
    \midrule
    Korpus kelompok & 10 & 750 & 116,49 & 679 & 71 \\
    Korpus G04 & 11 & 1.435 & 95,22 & 1.355 & 80 \\
    \midrule
    Total & 21 & 2.185 & 211,71 & 2.034 & 151 \\
    \bottomrule
  \end{tabular}
\end{table}

\section{Kendali mutu}

Kendali mutu Fase 1 memeriksa kejelasan, \textit{clipping}, derau, tumpang tindih ucapan, batas segmen, dan format audio. Pemeriksaan dilakukan melalui penyimakan, bentuk gelombang, serta spektrogram. Seleksi Fase 2 menerima label \textit{Good} dan \textit{Acceptable} serta membuang 21 ujaran G04 berlabel \textit{Reject}. Dua ujaran bertanda \kode{[unclear]} tetap digunakan dan ditandai untuk ditinjau \cite{fase1,korpusg04,repoasr}.

\section{Pembagian berdasarkan penutur}

Susunan 21 ID penutur diacak dengan \kode{numpy.random.RandomState(42)} setelah ID diurutkan. Sebanyak 13 ID pertama masuk ke himpunan latih, empat berikutnya ke validasi, dan empat terakhir ke uji. Pemeriksaan saat ekspor memastikan ID penutur dan berkas ujaran saling terpisah \cite{repoasr}.

\begin{table}[!htbp]
  \centering\small
  \caption{Pembagian data dan komposisi sumber.}
  \label{tab:split}
  \begin{tabular}{@{}lrrrrr@{}}
    \toprule
    Himpunan & Penutur & Ujaran & Durasi (menit) & Korpus kelompok & G04 \\
    \midrule
    Latih & 13 & 1.416 & 127,64 & 445 & 971 \\
    Validasi & 4 & 442 & 48,67 & 175 & 267 \\
    Uji & 4 & 327 & 35,40 & 130 & 197 \\
    \bottomrule
  \end{tabular}
\end{table}

Validasi memakai \kode{G04_S05}, \kode{G04_S08}, \kode{S02}, dan \kode{S10}. Uji memakai \kode{G04_S07}, \kode{G04_S11}, \kode{S04}, dan \kode{S09}. Daftar lengkap ada di \kode{ASR/data/splits/split_metadata.json}. Pembagian tunggal ini tidak diseimbangkan menurut gender, status penutur jati, atau derau.

\chapter{Pemodelan ASR}
\label{bab:pemodelan}

\section{Pemilihan teknik}

Kedua model memakai wav2vec 2.0 XLSR-53 dengan keluaran karakter. Pada prapelatihan wav2vec 2.0, sebagian representasi audio laten ditutupi. Model dilatih untuk membedakan target terkuantisasi yang sesuai dari kandidat lain. XLSR memperluas pendekatan ini ke audio lintas bahasa dengan representasi laten bersama \cite{baevski2020,conneau2021}. Dalam tugas ini, penyetelan menggunakan pasangan audio dan transkripsi Sunda dimulai dari model yang sudah memiliki kepala ASR. Pendekatan ini memanfaatkan representasi akustik dari tahap prapelatihan. Dengan demikian, model dapat diuji pada korpus lokal yang memiliki sekitar 128 menit audio latih.

\section{Arsitektur dan fungsi objektif}

Ekstraktor fitur tujuh lapis konvolusi mengubah sampel gelombang menjadi urutan vektor. Setiap lapis menghasilkan 512 kanal, dengan ukuran kernel $(10,3,3,3,3,2,2)$ dan langkah $(5,2,2,2,2,2,2)$. Langkah total 320 sampel menghasilkan sekitar satu bingkai setiap 20 milidetik pada 16 kHz. Setelah proyeksi fitur, enkoder Transformer dengan 24 lapis, dimensi tersembunyi 1.024, dan 16 kepala atensi per lapis membentuk representasi kontekstual. Lapisan linear menghasilkan logit per bingkai untuk setiap token CTC \cite{cahyamodel,nlpmodel,transformers}.

\begin{figure}[htbp]
  \centering
  \small
  \setlength{\fboxsep}{5pt}
  \fbox{Gelombang 16 kHz}
  \(\rightarrow\)
  \fbox{CNN}
  \(\rightarrow\)
  \fbox{Transformer}
  \(\rightarrow\)
  \fbox{Kepala CTC}
  \(\rightarrow\)
  \fbox{Dekode}
  \(\rightarrow\)
  \fbox{Teks}
  \caption{Alur inferensi wav2vec 2.0 untuk satu ujaran.}
  \label{fig:alur-asr}
\end{figure}

\textit{Connectionist Temporal Classification} (CTC) melatih model tanpa penyelarasan waktu pada tingkat karakter. Untuk audio \(x\) dan transkripsi \(y\), CTC menjumlahkan probabilitas semua jalur bingkai \(\pi\) yang dipetakan menjadi \(y\). Pemetaan \(B\) menggabungkan token identik yang berurutan, lalu menghapus token kosong (\textit{blank}). Token kosong di antara dua karakter identik mempertahankan kedua karakter pada keluaran \cite{graves2006}.
\begin{equation}
  \mathcal{L}_{\mathrm{CTC}}(x,y)
    = -\log \sum_{\pi\,:\,B(\pi)=y}
      \prod_{t=1}^{T} p(\pi_t \mid x).
\end{equation}

\clearpage
\section{Model pralatih dan kosakata}

\begin{description}
  \item[Cahya.] \kode{cahya/wav2vec2-large-xlsr-sundanese}, revisi \kode{beffe9e905ecd0787c3e87f271f8df7142f23b5b}. Model ini telah disetel untuk bahasa Sunda. Kepala CTC bawaannya memiliki 29 keluaran dan kosakatanya belum memuat \textit{é}. Sebelum penyetelan pada korpus gabungan, token \textit{é} dan satu baris baru pada kepala CTC ditambahkan. Bobot keluaran yang sudah ada disalin, sedangkan baris tambahan diinisialisasi secara acak \cite{cahyamodel,cahyaresep,repoasr}.
  \item[Indonesian NLP.] \kode{indonesian-nlp/wav2vec2-indonesian-javanese-sundanese}, revisi \kode{e5e699fa5aa5bdce999276a90a29589587b58ac9}. Model ini telah disetel menggunakan data bahasa Indonesia, Jawa, dan Sunda. Kosakata bawaannya sudah memuat \textit{é} dan kepala CTC memiliki 30 keluaran. Oleh karena itu, model ini dapat disetel langsung tanpa perluasan kosakata \cite{nlpmodel,repoasr}.
\end{description}

Kedua model hasil penyetelan memiliki 30 keluaran CTC. Token \kode{|} menandai batas kata. Token pengisi \kode{[PAD]} juga berfungsi sebagai \textit{blank} CTC. Token \kode{[UNK]} menandai karakter di luar kosakata.

\section{Konfigurasi penyetelan}

\begin{table}[htbp]
  \centering\small
  \caption{Konfigurasi pelatihan penuh kedua model.}
  \label{tab:konfigurasi}
  \begin{tabularx}{\textwidth}{@{}lY@{}}
    \toprule
    Komponen & Konfigurasi \\
    \midrule
    Data dan \textit{seed} & 1.416 ujaran latih, 442 ujaran validasi, \textit{seed} 42 \\
    Ukuran \textit{batch} & 4 per GPU, 2 GPU, akumulasi 4 langkah, \textit{batch} efektif 32 \\
    Laju pembelajaran & \(3 \times 10^{-4}\), pemanasan \pct{10}, peluruhan bobot 0 \\
    Durasi pelatihan & Maksimum 50 \textit{epoch}, penghentian dini setelah 5 evaluasi tanpa perbaikan WER \\
    Parameter dan memori & Ekstraktor konvolusi dibekukan, pelatihan FP16, \textit{gradient checkpointing} \\
    Regularisasi & \textit{Attention dropout} dan \textit{hidden dropout} 0,1, \textit{feature projection dropout} 0,0, \textit{layerdrop} 0,1, probabilitas \textit{time masking} 0,3 \\
    CTC dan pemilihan model & Reduksi \kode{mean}, \kode{ctc_zero_infinity} aktif, evaluasi dan penyimpanan per \textit{epoch}, WER validasi minimum \\
    \bottomrule
  \end{tabularx}
\end{table}

Konfigurasi pada Tabel~\ref{tab:konfigurasi} berasal dari \textit{notebook} penyetelan dan \kode{training_config.json} masing-masing model \cite{repoasr}. Enkoder Transformer dan kepala CTC tetap dilatih. Pembekuan hanya diterapkan pada ekstraktor fitur konvolusi.

\chapter{Implementasi dan Prapemrosesan}
\label{bab:implementasi}

\section{Penyusunan manifes dan validasi audio}

Skrip \kode{prepare_asr_dataset.py} menghubungkan metadata dan anotasi melalui nama berkas. Untuk G04, skrip menyelaraskan 45 nama berkas dalam anotasi yang kode topiknya berbeda dari metadata. Skrip lalu menambahkan awalan \kode{G04_} pada ID penutur dan ID ujaran untuk mencegah benturan dengan ID korpus kelompok. Skrip membaca durasi langsung dari audio. Sebanyak 64 selisih durasi di atas 0,2 detik dicatat dalam metadata pembagian \cite{repoasr}.

Skrip memeriksa keberadaan dan isi audio, format WAV mono PCM 16-bit, laju sampel 16 atau 48 kHz, label kualitas, dan kelengkapan transkripsi. Skrip juga menguji kelayakan panjang target CTC dengan menaksir jumlah bingkai dari durasi audio. Taksiran ini turut memperhitungkan token kosong tambahan untuk memisahkan karakter berulang \cite{graves2006,repoasr}.

\section{Prapemrosesan audio dan transkripsi}

Skrip \kode{export_kaggle_dataset.py} menyalin audio 16 kHz dan menurunkan laju sampel 245 berkas G04 dari 48 ke 16 kHz dengan filter polifase \kode{scipy.signal.resample_poly}. Hasilnya disimpan kembali sebagai PCM 16-bit. Pada model, ekstraktor fitur dari \kode{Wav2Vec2Processor} menormalisasi gelombang sebelum konvolusi. Masukan utama berupa gelombang satu kanal pada 16 kHz \cite{scipy,soundfile,transformers,repoasr}.

Normalisasi transkripsi versi \kode{v1} diterapkan sebelum pembagian data dan menghasilkan target yang sama untuk kedua model \cite{repoasr}.

\begin{enumerate}
  \item Ubah teks ke bentuk Unicode NFC dan huruf kecil. Pertahankan huruf \textit{é}.
  \item Hapus penanda \kode{[ee]}, \kode{[ee..]}, \kode{[noise]}, \kode{[overlap]}, dan \kode{[s]}. Buang penanda \kode{[unclear]} sambil menandai rekamannya untuk ditinjau. Pertahankan isi penanda lain tanpa kurung siku. Pertahankan pula kata pengisi tanpa penanda yang masih tercatat sebagai kata.
  \item Ubah angka 0 sampai 99 menjadi bentuk lisan bahasa Sunda. Pemetaan eksplisit juga mencakup \kode{s2} menjadi \textit{es dua} dan \kode{2023} menjadi \textit{dua rebu dua puluh tilu}.
  \item Perlakukan tanda hubung dan tanda baca sebagai pemisah kata. Rapatkan spasi berlebih dan pastikan seluruh karakter target termasuk dalam himpunan huruf yang didukung.
\end{enumerate}

\clearpage
\section{Pelatihan dan inferensi}

\textit{Notebook} \kode{01_pretrained_baseline.ipynb} menjalankan inferensi awal. \textit{Notebook} \kode{02_finetune_cahya.ipynb} dan \kode{03_finetune_indonesian_nlp.ipynb} memuat data ekspor Kaggle, memeriksa pembagian penutur, dan menyiapkan target karakter. Pada setiap \textit{batch}, panjang gelombang disamakan melalui penambahan \textit{padding}. Posisi \textit{padding} pada target diisi \kode{-100} agar diabaikan dalam perhitungan kerugian. Objek \kode{Trainer} menjalankan penyetelan dan menyimpan model sesuai WER validasi. Resep pelatihan mengacu pada \textit{notebook} Cahya untuk regularisasi model dan pembekuan ekstraktor fitur. Konfigurasi tersebut kemudian disesuaikan dengan pembagian penutur dan normalisasi korpus gabungan \cite{cahyaresep,transformers,accelerate,repoasr}.

Saat inferensi, gelombang dikelompokkan menurut panjang dan diproses dalam \textit{batch} berisi paling banyak empat ujaran. Dekode \textit{greedy} memilih token dengan logit terbesar pada setiap bingkai, menerapkan aturan penggabungan CTC, lalu memulihkan batas kata. Konfigurasi ini menggunakan keluaran model akustik tanpa model bahasa eksternal. \textit{Notebook} 03 menyertakan demonstrasi transkripsi menggunakan model yang disimpan oleh \textit{notebook} tersebut \cite{repoasr}.

\section{Lingkungan dan artefak}

Kedua pelatihan menggunakan Kaggle dan dua GPU Tesla T4. Versi yang tercatat mencakup Python 3.13.16, PyTorch 2.11.0 dengan CUDA 13.0, dan Transformers 5.20.2. \textit{Notebook} 03 mencatat pula Accelerate 1.15.0, JiWER 4.0.0, SoundFile 0.14.0, SciPy 1.16.3, pandas 2.3.3, dan NumPy 2.1.3. Matplotlib dipakai untuk kurva pelatihan \cite{pytorch,transformers,accelerate,jiwer,soundfile,scipy,pandas,numpy,matplotlib,repoasr}.

Repositori menyediakan konfigurasi, \textit{tokenizer}, riwayat pelatihan, dan hasil evaluasi pada \kode{ASR/results/}. \textit{Notebook} menyimpan bobot model dan prosesor sebagai keluaran Kaggle di bawah \kode{/kaggle/working/asr_outputs/}. Tautan unduhan bobot hasil penyetelan belum tercatat pada \textit{snapshot} repositori yang dirujuk.

\chapter{Skenario Eksperimen}
\label{bab:eksperimen}

\section{Model awal dan skenario penyetelan}
\textit{Bagian ini disiapkan untuk penulis Bab 5.}

\section{Protokol evaluasi}
\textit{Bagian ini disiapkan untuk penulis Bab 5.}

\chapter{Hasil dan Evaluasi}
\label{bab:hasil}

\section{Hasil validasi dan uji}
\textit{Bagian ini disiapkan untuk penulis Bab 6.}

\section{Contoh pengenalan ucapan}
\textit{Bagian ini disiapkan untuk penulis Bab 6.}

\chapter{Analisis Kesalahan}
\label{bab:analisis}

\section{Pola kesalahan dan kosakata}

Rincian operasi edit tingkat kata dan karakter pada model Indonesian NLP dianalisis untuk mengidentifikasi sumber utama galat. Tabel~\ref{tab:komposisi-edit} memuat jumlah substitusi, penghapusan, dan penyisipan pada himpunan validasi serta himpunan uji sebelum dan sesudah penyetelan \cite{repoasr}.

\begin{table}[!htbp]
  \centering\small
  \caption{Komposisi operasi edit kata dan karakter model Indonesian NLP pada himpunan validasi dan uji.}
  \label{tab:komposisi-edit}
  \begin{tabular}{@{}llrrrrrrrr@{}}
    \toprule
    & & \multicolumn{4}{c}{Tingkat kata} & \multicolumn{4}{c}{Tingkat karakter} \\
    \cmidrule(lr){3-6}\cmidrule(lr){7-10}
    Himpunan & Kondisi & Rujukan & Sub. & Del. & Ins. & Rujukan & Sub. & Del. & Ins. \\
    \midrule
    Validasi & \textit{Zero-shot} & 4.746 & 2.573 & 511 & 284 & 28.958 & 3.518 & 2.653 & 2.153 \\
    Validasi & Disetel (\textit{epoch} 14) & 4.746 & 1.383 & 308 & 283 & 28.958 & 2.462 & 1.960 & 1.489 \\
    \midrule
    Uji & \textit{Zero-shot} & 4.179 & 2.147 & 523 & 112 & 24.258 & 2.063 & 2.061 & 1.446 \\
    Uji & Disetel (\textit{epoch} 14) & 4.179 & 1.334 & 215 & 167 & 24.258 & 1.272 & 1.250 & 1.052 \\
    \bottomrule
  \end{tabular}
\end{table}

Pada himpunan uji setelah penyetelan, terdapat 1.716 operasi edit kata terhadap 4.179 kata rujukan. Substitusi mendominasi galat kata sebesar \pct{77,74} (1.334 kata), diikuti penghapusan sebesar \pct{12,53} (215 kata) dan penyisipan sebesar \pct{9,73} (167 kata). Dibandingkan kondisi \textit{zero-shot}, penyetelan memangkas substitusi kata uji sebesar \pct{37,87} dan penghapusan kata sebesar \pct{58,89}. Penyisipan kata sedikit bertambah dari 112 menjadi 167 karena model lebih aktif memecah segmen suara menjadi token kata.

Pada tingkat karakter uji, jumlah substitusi (1.272 karakter) dan penghapusan (1.250 karakter) hampir seimbang, disertai 1.052 penyisipan dari total 24.258 karakter rujukan. Perbedaan satu huruf saja pada dekode CTC tingkat karakter sudah cukup untuk mengubah satu kata menjadi substitusi kata penuh. Hal tersebut menjelaskan mengapa nilai WER uji (\pct{41,06}) hampir tiga kali lipat nilai CER uji (\pct{14,73}).

Tabel~\ref{tab:kesalahan-kata} menampilkan pasangan substitusi, kata terhapus, dan kata tersisip yang paling sering muncul pada himpunan uji \cite{repoasr}. Pola galat kata tersebut terbagi ke dalam empat fenomena linguistik dan ortografis utama.

\begin{table}[!htbp]
  \centering\small
  \caption{Kesalahan kata yang paling sering terjadi pada himpunan uji untuk model Indonesian NLP hasil penyetelan.}
  \label{tab:kesalahan-kata}
  \begin{tabular}{@{}lrlrlr@{}}
    \toprule
    Substitusi (\(\text{rujukan} \rightarrow \text{prediksi}\)) & Jumlah & Kata terhapus & Jumlah & Kata tersisip & Jumlah \\
    \midrule
    \(\textit{abi} \rightarrow \textit{abdi}\) & 39 & \textit{ee} & 29 & \textit{ka} & 8 \\
    \(\textit{ee} \rightarrow \textit{e}\) & 15 & \textit{nya} & 8 & \textit{a} & 4 \\
    \(\textit{ee} \rightarrow \textit{a}\) & 15 & \textit{the} & 5 & \textit{di} & 4 \\
    \(\textit{gitu} \rightarrow \textit{kitu}\) & 11 & \textit{di} & 4 & \textit{nu} & 4 \\
    \(\textit{téh} \rightarrow \textit{teh}\) & 10 & \textit{ka} & 4 & \textit{nya} & 4 \\
    \(\textit{unggal} \rightarrow \textit{ungal}\) & 9 & \textit{kitu} & 4 & \textit{sareng} & 3 \\
    \(\textit{anjeunna} \rightarrow \textit{anjeuna}\) & 7 & \textit{mah} & 4 & \textit{e} & 3 \\
    \(\textit{teh} \rightarrow \textit{téh}\) & 5 & \textit{in} & 4 & \textit{asa} & 2 \\
    \(\textit{minggu} \rightarrow \textit{mingu}\) & 5 & \textit{heumm} & 4 & \textit{ke} & 2 \\
    \(\textit{pun} \rightarrow \textit{mun}\) & 4 & \textit{nu} & 3 & \textit{téh} & 2 \\
    \(\textit{sanggeus} \rightarrow \textit{sangeus}\) & 4 & \textit{lah} & 3 & \textit{anu} & 2 \\
    \(\textit{kanggo} \rightarrow \textit{kango}\) & 4 & \textit{abdi} & 3 & \textit{mun} & 2 \\
    \bottomrule
  \end{tabular}
\end{table}

Pola pertama adalah variasi antara bentuk lisan dan ejaan baku bahasa Sunda. Substitusi terbanyak pada himpunan uji adalah \(\textit{abi} \rightarrow \textit{abdi}\) sebanyak 39 kali karena penutur uji melafalkan bentuk lisan \textit{abi}, sedangkan model cenderung menulis bentuk baku \textit{abdi} yang dominan di data latih. Pola serupa tampak pada pasangan alih kode \(\textit{gitu} \rightarrow \textit{kitu}\) (11 kali) dan bentuk singkat \(\textit{kusabab} \rightarrow \textit{sabab}\) (3 kali).

Pola kedua adalah peluruhan konsonan ganda pada kata seperti \(\textit{unggal} \rightarrow \textit{ungal}\) (9 kali), \(\textit{anjeunna} \rightarrow \textit{anjeuna}\) (7 kali), \(\textit{minggu} \rightarrow \textit{mingu}\) (5 kali), serta \(\textit{sanggeus} \rightarrow \textit{sangeus}\) (4 kali). Dalam aturan penggabungan CTC, dua huruf identik yang berurutan wajib dipisahkan oleh token kosong (\textit{blank}) agar tidak melebur menjadi satu huruf. Ketika durasi pelafalan konsonan ganda berlangsung singkat, model tidak menyisipkan token kosong di antaranya. Pola serupa juga mendominasi himpunan validasi melalui substitusi \(\textit{unggal} \rightarrow \textit{ungal}\) (12 kali) dan \(\textit{nganggo} \rightarrow \textit{ngango}\) (8 kali).

Pola ketiga adalah pertukaran dua arah antara huruf \textit{é} dan \textit{e}, terutama pada partikel \(\textit{téh} \rightarrow \textit{teh}\) (10 kali) dan \(\textit{teh} \rightarrow \textit{téh}\) (5 kali). Inkonsistensi penulisan \textit{teh} dan \textit{téh} pada sebagian anotasi rujukan membuat kedua ejaan tersebut saling bersaing. Meskipun demikian, selisih antara WER ketat (\pct{41,06}) dan WER tanpa aksen (\pct{39,91}) pada himpunan uji hanya 1,15 poin persentase.

Pola keempat berkaitan dengan kata pengisi, partikel pendek, dan anomali anotasi. Pada korpus G04, gumaman \textit{ee} dan \textit{heumm} ditulis langsung sebagai kata biasa tanpa tanda kurung siku sehingga tidak terhapus oleh normalisasi \kode{v1}. Akibatnya, token \textit{ee} menjadi kata yang paling sering terhapus (29 kali) atau berubah menjadi \textit{e} dan \textit{a} (masing-masing 15 kali). Selain itu, kata bahasa Inggris \textit{the} (5 kali) dan \textit{in} (4 kali) muncul pada daftar kata terhapus karena beberapa rujukan penutur \kode{G04_S11} berisi kalimat bahasa Inggris meskipun audionya berbahasa Sunda.

Pengaruh cakupan kosakata latih terhadap tingkat kesalahan kata disajikan pada Tabel~\ref{tab:kata-jarang} \cite{repoasr}. Himpunan latih memuat 2.201 kata unik. Setiap kata rujukan pada himpunan validasi dan uji dikelompokkan berdasarkan keberadaannya di dalam kosakata latih tersebut.

\begin{table}[!htbp]
  \centering\small
  \caption{Tingkat kesalahan kata berdasarkan ketercakupan dalam kosakata data latih (2.201 kata unik).}
  \label{tab:kata-jarang}
  \begin{tabular}{@{}llrrr@{}}
    \toprule
    Himpunan & Kelompok kosakata & Jumlah kata & Kata salah & Persentase salah (\pct{}) \\
    \midrule
    Validasi & Ada di data latih & 4.183 & 1.345 & 32,15 \\
    Validasi & Tidak ada di data latih (\textit{out-of-vocabulary}) & 563 & 346 & 61,46 \\
    \midrule
    Uji & Ada di data latih & 3.456 & 1.091 & 31,57 \\
    Uji & Tidak ada di data latih (\textit{out-of-vocabulary}) & 723 & 458 & 63,35 \\
    \bottomrule
  \end{tabular}
\end{table}

Pada himpunan uji, sebanyak 723 dari 4.179 kata rujukan (\pct{17,30}) tidak pernah muncul dalam data latih. Proporsi kesalahan pada kata di luar kosakata latih tersebut mencapai \pct{63,35} (458 kata), atau dua kali lipat tingkat kesalahan pada kata yang sudah ada di data latih (\pct{31,57}). Perbandingan serupa terjadi pada himpunan validasi dengan tingkat kesalahan \pct{61,46} berbanding \pct{32,15}. Karena dekode \textit{greedy} CTC menyusun huruf tanpa model bahasa, nama diri, kata serapan, dan kata berimbuhan baru mudah mengalami selisih satu atau dua karakter.

\section{Variasi antarpenutur dan kondisi rekaman}

Tingkat kesalahan juga bervariasi tajam menurut karakteristik penutur dan rekaman. Tabel~\ref{tab:kelompok-uji} merangkum WER dan CER beserta ukuran sampel (\(n\)) untuk setiap kelompok penutur, status penutur jati, label kualitas rekaman, dan rentang durasi pada himpunan validasi serta himpunan uji \cite{repoasr}.

\begin{table}[!htbp]
  \centering\small
  \caption{Rincian WER dan CER model Indonesian NLP menurut kelompok penutur, status penutur jati, kualitas audio, dan durasi ujaran.}
  \label{tab:kelompok-uji}
  \begin{tabular}{@{}llrrrrrrr@{}}
    \toprule
    & \multicolumn{4}{c}{Himpunan validasi (\(n = 442\))} & \multicolumn{4}{c}{Himpunan uji (\(n = 327\))} \\
    \cmidrule(lr){2-5}\cmidrule(lr){6-9}
    Kelompok & Nilai & \(n\) & WER (\pct{}) & CER (\pct{}) & Nilai & \(n\) & WER (\pct{}) & CER (\pct{}) \\
    \midrule
    Penutur & \kode{S02} & 69 & 13,66 & 2,52 & \kode{S04} & 69 & 19,73 & 4,28 \\
    Penutur & \kode{G04_S05} & 84 & 31,67 & 10,37 & \kode{S09} & 61 & 24,80 & 4,98 \\
    Penutur & \kode{S10} & 106 & 35,22 & 9,28 & \kode{G04_S07} & 115 & 59,84 & 25,35 \\
    Penutur & \kode{G04_S08} & 183 & 82,26 & 58,83 & \kode{G04_S11} & 82 & 75,98 & 35,83 \\
    \midrule
    Penutur jati & \textit{Non-native} & 69 & 13,66 & 2,52 & \textit{Native} & 130 & 21,88 & 4,59 \\
    Penutur jati & \textit{Native} & 106 & 35,22 & 9,28 & \textit{Unknown} & 197 & 64,40 & 28,28 \\
    Penutur jati & \textit{Unknown} & 267 & 57,67 & 35,17 & & & & \\
    \midrule
    Kualitas & \textit{Good} & 440 & 41,58 & 20,43 & \textit{Good} & 278 & 37,76 & 12,67 \\
    Kualitas & \textit{Acceptable} & 2 & 44,44 & 15,97 & \textit{Acceptable} & 49 & 62,43 & 29,01 \\
    \midrule
    Durasi & \(< 3\text{ s}\) & 79 & 89,37 & 65,49 & \(< 3\text{ s}\) & 1 & 62,50 & 37,84 \\
    Durasi & \(3\text{ sampai }6\text{ s}\) & 138 & 67,08 & 43,82 & \(3\text{ sampai }6\text{ s}\) & 212 & 60,72 & 26,04 \\
    Durasi & \(6\text{ sampai }10\text{ s}\) & 140 & 28,82 & 8,63 & \(6\text{ sampai }10\text{ s}\) & 79 & 19,82 & 4,30 \\
    Durasi & \(> 10\text{ s}\) & 85 & 28,25 & 6,66 & \(> 10\text{ s}\) & 35 & 25,58 & 5,34 \\
    \bottomrule
  \end{tabular}
\end{table}

Pada himpunan uji, terdapat kesenjangan mencolok antara penutur korpus kelompok dan penutur korpus G04. Dua penutur jati dari korpus kelompok, yakni \kode{S04} (\(n = 69\)) dan \kode{S09} (\(n = 61\)), masing-masing mencapai WER \pct{19,73} (CER \pct{4,28}) dan WER \pct{24,80} (CER \pct{4,98}). Sebaliknya, penutur \kode{G04_S07} (\(n = 115\)) mencatat WER \pct{59,84} dan \kode{G04_S11} (\(n = 82\)) mencatat WER \pct{75,98}. Pola serupa terjadi pada himpunan validasi ketika penutur \kode{G04_S08} (\(n = 183\)) mencatat WER \pct{82,26} dan CER \pct{58,83} akibat sejumlah segmen transkripsi yang tidak selaras dengan potongan audionya.

Perbedaan berdasarkan status penutur jati berimpit langsung dengan asal korpus dan gaya bicara. Pada himpunan uji, kelompok \textit{Native} (\(n = 130\), gabungan \kode{S04} dan \kode{S09}) memperoleh WER \pct{21,88} dan CER \pct{4,59}, sedangkan kelompok \textit{Unknown} (\(n = 197\), seluruhnya dari G04) memperoleh WER \pct{64,40} dan CER \pct{28,28}. Pada himpunan validasi, penutur \textit{Non-native} \kode{S02} (\(n = 69\)) justru meraih WER terendah (\pct{13,66}) karena penutur tersebut membacakan naskah dengan artikulasi teratur, sementara penutur \textit{Native} \kode{S10} (\(n = 106\)) berbicara secara spontan dengan WER \pct{35,22}.

Ditinjau dari mutu rekaman pada himpunan uji, 278 ujaran berlabel \textit{Good} menghasilkan WER \pct{37,76} dan CER \pct{12,67}. Sebaliknya, 49 ujaran berlabel \textit{Acceptable} yang memiliki derau latar atau artikulasi kurang jelas menghasilkan WER \pct{62,43} dan CER \pct{29,01}. Pada himpunan validasi, sampel berlabel \textit{Acceptable} hanya berjumlah dua ujaran (\(n = 2\)) sehingga perbandingannya terhadap label \textit{Good} (\(n = 440\)) tidak representatif secara statistik.

Kelompok durasi memperlihatkan bahwa ujaran berdurasi 6 sampai 10 detik pada himpunan uji (\(n = 79\), WER \pct{19,82}) dan di atas 10 detik (\(n = 35\), WER \pct{25,58}) jauh lebih akurat daripada ujaran 3 sampai 6 detik (\(n = 212\), WER \pct{60,72}). Pola ini tidak semata-mata disebabkan oleh panjang konteks akustik, melainkan berbaur dengan sumber korpus. Rekaman korpus G04 dipotong menjadi segmen pendek berdurasi di bawah enam detik, sedangkan korpus kelompok memuat kalimat utuh yang lebih panjang dan bersih.

\chapter{Penutup}
\label{bab:penutup}

\section{Kesimpulan}

Penyetelan ulang model pralatih wav2vec 2.0 XLSR-53 berbasis CTC menggunakan 1.416 ujaran latih (127,64 menit audio dari 13 penutur) berhasil membangun sistem ASR bahasa Sunda untuk ujaran spontan. Penyesuaian normalisasi transkripsi \kode{v1} dan perluasan kepala CTC pada model Cahya memungkinkan kedua model mempertahankan perbedaan ortografis antara huruf \textit{e} dan \textit{é}. Hasil penapisan empat kandidat menunjukkan bahwa laju pembelajaran \(3 \times 10^{-4}\) memberikan konvergensi terbaik pada kedua model.

Pada himpunan validasi (442 ujaran dari empat penutur), penyetelan menurunkan WER model Cahya dari \pct{67,28} pada kondisi \textit{zero-shot} menjadi \pct{44,44} dan CER dari \pct{27,60} menjadi \pct{20,94}. Pada himpunan validasi yang sama, model Indonesian NLP mengalami penurunan WER yang lebih besar dari \pct{70,97} menjadi \pct{41,59} serta penurunan CER dari \pct{28,75} menjadi \pct{20,41}. Pada 327 ujaran uji dari empat penutur terpisah, model Indonesian NLP terpilih menurunkan WER dari \pct{66,57} menjadi \pct{41,06}, menurunkan CER dari \pct{22,96} menjadi \pct{14,73}, dan meningkatkan akurasi kata berhuruf \textit{é} dari \pct{0,00} menjadi \pct{74,65}.

Analisis kesalahan pada himpunan uji menunjukkan bahwa substitusi kata menyumbang \pct{77,74} dari seluruh operasi edit kata, sedangkan kesalahan tingkat karakter jauh lebih rendah dibanding tingkat kata. Kesalahan kata terutama dipicu oleh kata di luar kosakata latih yang tingkat salahnya mencapai \pct{63,35}, dibandingkan \pct{31,57} pada kata yang tercakup di data latih. Sumber kesalahan utama lainnya meliputi perbedaan bentuk lisan terhadap ejaan baku (\textit{abi} dan \textit{abdi}), peluruhan konsonan ganda pada dekode CTC, kata pengisi \textit{ee} tanpa penanda kurung siku, serta anomali penyelarasan anotasi pada sebagian penutur korpus G04.

\section{Keterbatasan}

Eksperimen ini memiliki beberapa keterbatasan metodologis dan data yang perlu diperhatikan dalam menafsirkan hasil:

\begin{enumerate}
  \item Evaluasi menggunakan satu pembagian acak berbasis penutur dengan empat penutur validasi dan empat penutur uji tanpa penyeimbangan menurut gender, tingkat derau, maupun gaya bicara baca terhadap spontan.
  \item Seluruh sebelas penutur korpus G04 memiliki status penutur jati \textit{Unknown}, dan sebagian berkas pada penutur seperti \kode{G04_S08} serta \kode{G04_S11} memuat ketidakselarasan batas potongan audio atau teks rujukan berbahasa Inggris yang meningkatkan angka WER secara semu.
  \item Penapisan laju pembelajaran hanya dijalankan selama lima \textit{epoch} dengan satu \textit{seed} acak, sehingga laju pembelajaran yang lebih kecil berpotensi tampak kurang kompetitif karena pembaruan bobotnya berlangsung lebih lambat.
  \item Tahap inferensi masih mengandalkan dekode \textit{greedy} CTC tingkat karakter tanpa model bahasa eksternal, sehingga kesalahan satu karakter langsung menghasilkan kata tidak baku yang menurunkan WER.
\end{enumerate}

\section{Saran perbaikan}

Berdasarkan pola kesalahan dan keterbatasan tersebut, tiga arah perbaikan direkomendasikan untuk pengembangan berikutnya:

\begin{enumerate}
  \item \textbf{Pembersihan data dan pembakuan ortografi.} Lakukan audit ulang keselarasan audio dan teks pada korpus G04 untuk membuang atau memperbaiki rujukan yang terpotong serta rujukan berbahasa Inggris pada penutur \kode{G04_S08} dan \kode{G04_S11}. Seragamkan pula aturan penandaan kata pengisi seperti \textit{ee} dan \textit{heumm}, serta susun kamus pemetaan varian lisan ke bentuk baku seperti \textit{abi} menjadi \textit{abdi} dan \textit{gitu} menjadi \textit{kitu}.
  \item \textbf{Perbaikan protokol evaluasi.} Terapkan validasi silang lipat ganda berbasis penutur (\textit{speaker-grouped cross-validation}) dengan stratifikasi berdasarkan sumber korpus, kualitas rekaman, dan gaya bicara. Langkah ini mengurangi kepekaan metrik evaluasi terhadap karakteristik empat penutur tertentu pada satu pembagian tunggal.
  \item \textbf{Integrasi model bahasa dan eksplorasi arsitektur.} Gunakan dekode pencarian berkas (\textit{beam search}) CTC yang dipadukan dengan model bahasa \(n\)-gram bahasa Sunda untuk memperbaiki peluruhan konsonan ganda, batas kata, dan ejaan kata jarang. Selain itu, uji arsitektur pengenal ucapan multibahasa yang lebih baru seperti Whisper atau Massively Multilingual Speech (MMS) pada pembagian penutur yang identik sebagai pembanding.
\end{enumerate}


\clearpage
\phantomsection
\addcontentsline{toc}{chapter}{Daftar Pustaka}
\renewcommand{\bibname}{Daftar Pustaka}
\bibliographystyle{unsrt}
\bibliography{references}
\end{document}
